\documentclass[12pt,a4paper]{article}
\usepackage[T2A]{fontenc}
\usepackage[utf8]{inputenc}
\usepackage[russian]{babel}
\usepackage{float}
\usepackage{amsmath, amssymb}
\usepackage{geometry}
\usepackage{tikz}
\usepackage{tikz-3dplot}
\usepackage{caption}
\usetikzlibrary{arrows.meta, calc, decorations.pathmorphing}
\geometry{margin=2cm}

\newcommand{\vn}{\bar n}
\newcommand{\vq}{\bar q}
\newcommand{\vr}{\bar r}
\newcommand{\vrzero}{\bar r_0}
\newcommand{\ve}{\bar e}
\newcommand{\vi}{\bar i}
\newcommand{\vj}{\bar j}
\newcommand{\vk}{\bar k}

\begin{document}

\setcounter{page}{58}

\section*{ЛЕКЦИЯ 9. ЛИНИИ НА ПЛОСКОСТИ. УРАВНЕНИЯ ПРЯМОЙ НА
ПЛОСКОСТИ. КРИВЫЕ 2-ГО ПОРЯДКА}

\subsection*{9.1. Уравнения прямой на плоскости}

Пусть на плоскости задана декартова система координат и некоторая линия $L$.

\begin{figure}[H]
  \centering
  \begin{minipage}{0.4\textwidth}
    \centering
    \begin{tikzpicture}[scale=0.9, >=Latex]
      \draw[->] (-0.2,0) -- (4.2,0) node[right] {$x$};
      \draw[->] (0,-0.2) -- (0,3.2) node[above] {$y$};
      \node[below left] at (0,0) {$0$};
      \draw[thick] (0.7,0.6) .. controls (1.4,2.4) and (2.8,2.8) .. (3.5,1.1);
      \node at (2.7,2.3) {$L$};
      \fill (2.1,2.05) circle (1.5pt);
      \node[right] at (2.1,2.05) {$M(x,y)$};
    \end{tikzpicture}
    \caption*{Рисунок 9.1}
  \end{minipage}%
  \begin{minipage}{0.55\textwidth}
    \textbf{Определение.} Уравнение $\Phi(x,y)=0$ называется
    уравнением линии $L$, если этому уравнению удовлетворяют
    координаты любой точки этой линии, и не удовлетворяют координаты
    точек, не принадлежащих линии $L$.
  \end{minipage}
\end{figure}

\subsection*{\textbf{Уравнение прямой на плоскости}}

\textbf{Теорема.} Любая прямая на плоскости задается уравнением вида
\begin{equation}
  Ax+By+C=0. \tag{9.1}
\end{equation}
И обратно: всякое уравнение вида (9.1) определяет на плоскости прямую.

\subsection*{\textbf{Доказательство}}

$\vn$ --- нормальный вектор.

Возьмем произвольную точку $M(x,y)\in L$.

\begin{figure}[H]
  \centering
  \begin{minipage}{0.4\textwidth}
    \centering
    \begin{tikzpicture}[scale=0.9, >=Latex]
      \draw[->] (-0.2,0) -- (4.5,0) node[right] {$x$};
      \draw[->] (0,-0.2) -- (0,3.4) node[above] {$y$};
      \node[below left] at (0,0) {$0$};
      \draw[thick] (0.5,0.7) -- (4.0,2.0) node[right] {$L$};
      \coordinate (M0) at (1.4,1.03);
      \coordinate (M) at (3.0,1.62);
      \fill (M0) circle (1.5pt);
      \fill (M) circle (1.5pt);
      \node[below] at (M0) {$M_0(x_0,y_0)$};
      \node[above] at (M) {$M(x,y)$};
      \draw[->, thick] (M0) -- ++(-0.45,1.25) node[above] {$\vn=(A,B)$};
      \node[left] at (1.0,2.05) {$\vn\perp L$};
    \end{tikzpicture}
    \caption*{Рисунок 9.2}
  \end{minipage}%
  \begin{minipage}{0.55\textwidth}
    $$
    \text{т. }M(x,y)\in L
    \Leftrightarrow \vn\perp \overline{M_0M}
    \Leftrightarrow
    $$
    $$
    \Leftrightarrow (\vn,\overline{M_0M})=0
    \Leftrightarrow
    $$
    \begin{equation}
      \Leftrightarrow A(x-x_0)+B(y-y_0)=0. \tag{9.2}
    \end{equation}
  \end{minipage}
\end{figure}

Раскроем скобки:
$$
Ax-Ax_0+By-By_0=[-Ax_0-By_0=C]=0
\Rightarrow
$$
$$
\Rightarrow Ax+By+C=0.
$$

Проводя обратные рассуждения, получим, что уравнение вида (9.1)
определяет на плоскости прямую.

\textbf{ч.т.д.}

(9.1) --- общее уравнение прямой.

(9.2) --- уравнение прямой, проходящей через точку $M_0(x_0,y_0)$
перпендикулярно нормальному вектору $\vn=(A,B)$.

Пусть в (9.1) $B\ne 0$:
$$
By=-Ax-C \Rightarrow y=-\frac ABx-\frac CB.
$$

Обозначим
$$
k=-\frac AB;\qquad b=-\frac CB,
$$
получим уравнение
\begin{equation}
  y=kx+b. \tag{9.3}
\end{equation}

(9.3) --- уравнение прямой с угловым коэффициентом.

Нетрудно показать, что $k=\tg\alpha$; $b=OB$.

\begin{figure}[H]
  \centering
  \begin{minipage}{0.4\textwidth}
    \centering
    \begin{tikzpicture}[scale=0.95, >=Latex]
      \draw[->] (-0.2,0) -- (4.4,0) node[right] {$x$};
      \draw[->] (0,-0.2) -- (0,3.4) node[above] {$y$};
      \node[below left] at (0,0) {$0$};
      \draw[thick] (0,0.9) -- (4,2.8) node[right] {$L$};
      \fill (0,0.9) circle (1.3pt);
      \node[left] at (0,0.9) {$B$};
      \draw[dashed] (0,0) -- (0,0.9);
      \node[left] at (0,0.45) {$b$};
      \draw (1.35,0) arc (0:25:1.35);
      \node at (1.55,0.35) {$\alpha$};
    \end{tikzpicture}
    \caption*{Рисунок 9.3}
  \end{minipage}%
  \begin{minipage}{0.55\textwidth}
    Прямая $L$ пересекает ось $Oy$ в точке $B$, причем отрезок
    $OB=b$, а угол наклона прямой к положительному направлению оси
    $Ox$ равен $\alpha$.
  \end{minipage}
\end{figure}

Пусть в (9.1) $B=0$:
$$
Ax+C=0 \Rightarrow x=-\frac CA.
$$

Обозначим
$$
a=-\frac CA,
$$
получим уравнение
\begin{equation}
  x=a. \tag{9.4}
\end{equation}

(9.4) --- уравнение вертикальной прямой.

\begin{figure}[H]
  \centering
  \begin{minipage}{0.4\textwidth}
    \centering
    \begin{tikzpicture}[scale=0.95, >=Latex]
      \draw[->] (-0.2,0) -- (4.4,0) node[right] {$x$};
      \draw[->] (0,-0.2) -- (0,3.4) node[above] {$y$};
      \node[below left] at (0,0) {$0$};
      \draw[thick] (2.2,0) -- (2.2,3.1) node[above] {$L$};
      \node[below] at (2.2,0) {$a$};
    \end{tikzpicture}
    \caption*{Рисунок 9.4}
  \end{minipage}%
  \begin{minipage}{0.55\textwidth}
    Вертикальная прямая задается уравнением $x=a$.
  \end{minipage}
\end{figure}

Уравнение (9.1) можно записать в виде
$$
\frac{x}{-C/A}+\frac{y}{-C/B}=1.
$$

Обозначим
$$
a=-\frac CA;\qquad b=-\frac CB
\Rightarrow
$$
\begin{equation}
  \Rightarrow \frac xa+\frac yb=1. \tag{9.5}
\end{equation}

(9.5) --- уравнение прямой в отрезках.

\begin{figure}[H]
  \centering
  \begin{minipage}{0.4\textwidth}
    \centering
    \begin{tikzpicture}[scale=0.9, >=Latex]
      \draw[->] (-0.2,0) -- (4.4,0) node[right] {$x$};
      \draw[->] (0,-0.2) -- (0,3.5) node[above] {$y$};
      \node[below left] at (0,0) {$0$};
      \coordinate (A) at (3.4,0);
      \coordinate (B) at (0,2.7);
      \draw[thick] (B) -- (A) node[midway, above right] {$L$};
      \fill (A) circle (1.3pt);
      \fill (B) circle (1.3pt);
      \node[below] at (A) {$(a,0)$};
      \node[left] at (B) {$(0,b)$};
      \node[below] at (2.1,0) {$a$};
      \node[left] at (0,1.35) {$b$};
    \end{tikzpicture}
    \caption*{Рисунок 9.5}
  \end{minipage}%
  \begin{minipage}{0.55\textwidth}
    Прямая отсекает на осях координат отрезки $a$ и $b$.
  \end{minipage}
\end{figure}

\subsection*{\textbf{Параметрическое уравнение прямой}}

$\vq=(l,m)$ --- направляющий вектор.

$$
\overline{M_0M}=t\cdot \vq,\qquad t\in\mathbb R.
$$

$t$ --- некоторый параметр.

$$
\overline{M_0M}=\vr-\vrzero=t\cdot \vq
\Rightarrow
$$
\begin{equation}
  \Rightarrow \vr=\vrzero+t\cdot \vq. \tag{9.6}
\end{equation}

или в координатной форме
\begin{equation}
  \begin{cases}
    x=x_0+t\cdot l,\\
    y=y_0+t\cdot m.
  \end{cases}
  \tag{9.7}
\end{equation}

\begin{figure}[H]
  \centering
  \begin{minipage}{0.4\textwidth}
    \centering
    \begin{tikzpicture}[scale=0.85, >=Latex]
      \draw[->] (-0.2,0) -- (5.0,0) node[right] {$x$};
      \draw[->] (0,-0.2) -- (0,3.6) node[above] {$y$};
      \node[below left] at (0,0) {$0$};
      \coordinate (M0) at (1.5,1.0);
      \coordinate (M) at (3.5,2.2);
      \draw[thick] (0.8,0.58) -- (4.3,2.68) node[right] {$L$};
      \fill (M0) circle (1.3pt);
      \fill (M) circle (1.3pt);
      \node[below] at (M0) {$M_0(x_0,y_0)$};
      \node[above] at (M) {$M(x,y)$};
      \draw[->] (0,0) -- (M0) node[midway, above left] {$\vrzero$};
      \draw[->] (0,0) -- (M) node[midway, below right] {$\vr$};
      \draw[->, thick] (M0) -- (M) node[midway, above] {$\vq$};
      \node at (3.1,1.25) {$\vq\parallel L$};
    \end{tikzpicture}
    \caption*{Рисунок 9.6}
  \end{minipage}%
  \begin{minipage}{0.55\textwidth}
    (9.6) и (9.7) --- параметрические уравнения прямой.

    Из (9.7):
    \begin{equation}
      \Rightarrow \frac{x-x_0}{l}=\frac{y-y_0}{m}. \tag{9.8}
    \end{equation}

    (9.8) --- каноническое уравнение прямой, проходящей через т.
    $M_0(x_0,y_0)$ с направляющим вектором $\vq=(l,m)$.
  \end{minipage}
\end{figure}

\subsection*{\textbf{Расстояние от точки до прямой}}

\begin{figure}[H]
  \centering
  \begin{minipage}{0.4\textwidth}
    \centering
    \begin{tikzpicture}[scale=0.9, >=Latex]
      \draw[thick] (-0.2,0.5) -- (4.2,1.0);
      \coordinate (M) at (1.0,0.64);
      \coordinate (M0) at (1.8,2.6);
      \coordinate (P) at (1.63,0.71);
      \fill (M) circle (1.3pt);
      \fill (M0) circle (1.3pt);
      \node[below] at (M) {$M(x,y)$};
      \node[above] at (M0) {$M_0(x_0,y_0)$};
      \draw[thick] (M) -- (M0);
      \draw[dashed] (M0) -- (P);
      \node[right] at ($(M0)!0.5!(P)$) {$d$};
      \draw[->, thick] (2.4,0.8) -- (2.7,2.7);
      \node[right] at (2.7,1.75) {$\vn=(A,B)$};
      \draw ($(P)+(0.12,0.02)$) -- ++(-0.02,0.18) -- ++(-0.16,-0.02);
    \end{tikzpicture}
    \caption*{Рисунок 9.7}
  \end{minipage}%
  \begin{minipage}{0.55\textwidth}
    $$
    \overline{MM_0}=(x_0-x,y_0-y);
    $$
    $$
    d=\left|\text{пр}_{\vn}\overline{MM_0}\right|
    =
    \left|\frac{(\vn,\overline{MM_0})}{|\vn|}\right|
    =
    \left|\frac{A(x_0-x)+B(y_0-y)}{|\vn|}\right|
    =
    $$
    $$
    =
    \left|\frac{Ax_0+By_0+(-Ax-By)}{|\vn|}\right|
    =
    \left|\frac{Ax_0+By_0+C}{\sqrt{A^2+B^2}}\right|;
    $$
    \begin{equation}
      d=\left|\frac{Ax_0+By_0+C}{\sqrt{A^2+B^2}}\right|. \tag{9.9}
    \end{equation}
  \end{minipage}
\end{figure}

\subsection*{\textbf{Направляющие косинусы}}

\begin{figure}[H]
  \centering
  \begin{minipage}{0.4\textwidth}
    \centering
    \tdplotsetmaincoords{65}{115}
    \begin{tikzpicture}[scale=0.9, tdplot_main_coords, >=Latex]
      \draw[->] (0,0,0) -- (3,0,0) node[below left] {$x$};
      \draw[->] (0,0,0) -- (0,3,0) node[right] {$y$};
      \draw[->] (0,0,0) -- (0,0,3) node[above] {$z$};
      \node[below] at (0,0,0) {$0$};
      \draw[->, thick] (0,0,0) -- (1.8,1.6,2.2) node[above] {$\vr=(x,y,z)$};
      \node at (0.45,0,0) {$\vi$};
      \node at (0,0.45,0) {$\vj$};
      \node at (0,0,0.55) {$\vk$};
      \node at (0.35,0.25,0.4) {$e$};
    \end{tikzpicture}
    \caption*{Рисунок 9.8}
  \end{minipage}%
  \begin{minipage}{0.55\textwidth}
    Обозначим
    $$
    \alpha=(\hat{\vr},\vi);\qquad
    \beta=(\hat{\vr},\vj);
    $$
    $$
    \gamma=(\hat{\vr},\vk).
    $$

    $\cos\alpha$, $\cos\beta$ и $\cos\gamma$ --- направляющие
    косинусы вектора $\vr$.

    Т.к.
    $$
    x=\text{пр}_{\vi}\vr;\qquad
    y=\text{пр}_{\vj}\vr;
    $$
    $$
    z=\text{пр}_{\vk}\vr
    $$
    $$
    \Rightarrow x=|\vr|\cos\alpha;\qquad
    y=|\vr|\cos\beta;
    $$
    $$
    z=|\vr|\cos\gamma.
    $$
  \end{minipage}
\end{figure}

$$
\vr=x\vi+y\vj+z\vk
=
|\vr|(\cos\alpha\,\vi+\cos\beta\,\vj+\cos\gamma\,\vk).
$$

Для единичного вектора
$$
\ve=\frac{\vr}{|\vr|}
$$
получим
\begin{equation}
  \ve=(\cos\alpha,\cos\beta,\cos\gamma). \tag{9.10}
\end{equation}

Из (9.10) следует, что
$$
\cos^2\alpha+\cos^2\beta+\cos^2\gamma=1.
$$

\subsection*{\textbf{Нормальное уравнение прямой}}

$$
\ve=\frac{\vn}{|\vn|};
\qquad
\text{где }\ve=(\cos\alpha,\cos\beta);
$$
$$
\overline{OM}=(x,y).
\quad
\text{Обозначим }p=|\vn|.
$$

$$
p=\text{пр}_{\ve}\overline{OM}
=
\frac{(\overline{OM},\ve)}{|\ve|}
=
$$
$$
=x\cos\alpha+y\cos\beta
$$
\begin{equation}
  \Rightarrow x\cos\alpha+y\cos\beta-p=0. \tag{9.11}
\end{equation}

\begin{figure}[H]
  \centering
  \begin{minipage}{0.4\textwidth}
    \centering
    \begin{tikzpicture}[scale=0.95, >=Latex]
      \draw[->] (-0.2,0) -- (4.4,0) node[right] {$x$};
      \draw[->] (0,-0.2) -- (0,3.4) node[above] {$y$};
      \node[below left] at (0,0) {$0$};
      \draw[thick] (1.1,0.3) -- (3.8,2.9) node[right] {$L$};
      \coordinate (M) at (2.8,1.93);
      \fill (M) circle (1.3pt);
      \node[right] at (M) {$M(x,y)$};
      \draw[->, thick] (0,0) -- (1.95,1.45) node[midway, above left] {$\vn$};
      \node at (1.6,2.1) {$\vn\perp L$};
      \node at (1.15,0.95) {$e$};
    \end{tikzpicture}
    \caption*{Рисунок 9.9}
  \end{minipage}%
  \begin{minipage}{0.55\textwidth}
    (9.11) --- нормальное уравнение прямой, где $\cos\alpha$ и
    $\cos\beta$ --- направляющие косинусы нормального вектора,
    направленного из начала координат в сторону прямой, а $p$ ---
    расстояние от начала координат до прямой.
  \end{minipage}
\end{figure}

\subsection*{\textbf{Угол между прямыми}}

$$
L_1:A_1x+B_1y+C_1=0;
\qquad
L_2:A_2x+B_2y+C_2=0.
$$

Обозначим
$$
\alpha=(\hat L_1,L_2).
$$

\begin{equation}
  \cos\alpha=
  \left|
  \frac{A_1A_2+B_1B_2}
  {\sqrt{A_1^2+B_1^2}\cdot\sqrt{A_2^2+B_2^2}}
  \right|;
  \qquad
  \frac{A_1}{A_2}=\frac{B_1}{B_2}. \tag{9.12}
\end{equation}

\begin{equation}
  A_1\cdot A_2+B_1\cdot B_2=0. \tag{9.13}
\end{equation}

(9.12) --- условие параллельности прямых $L_1$ и $L_2$.

(9.13) --- условие перпендикулярности прямых $L_1$ и $L_2$.

\subsection*{9.2. Кривые второго порядка}

\textbf{Определение.} Окружность --- это геометрическое место точек
на плоскости, равноудаленных от данной точки, называемой центром.

$$
R=|\overline{M_0M}|=\sqrt{(x-x_0)^2+(y-y_0)^2}
$$
\begin{equation}
  \Rightarrow (x-x_0)^2+(y-y_0)^2=R^2. \tag{9.14}
\end{equation}

(9.14) --- уравнение окружности с центром в т. $M_0(x_0,y_0)$ и радиусом $R$.

\begin{figure}[H]
  \centering
  \begin{minipage}{0.4\textwidth}
    \centering
    \begin{tikzpicture}[scale=0.9, >=Latex]
      \draw[->] (-0.2,0) -- (4.4,0) node[right] {$x$};
      \draw[->] (0,-0.2) -- (0,3.7) node[above] {$y$};
      \node[below left] at (0,0) {$0$};
      \coordinate (O) at (2.1,1.8);
      \draw[thick] (O) circle (1.15);
      \fill (O) circle (1.3pt);
      \node[left] at (O) {$M_0(x_0,y_0)$};
      \coordinate (M) at ($(O)+(45:1.15)$);
      \fill (M) circle (1.3pt);
      \node[right] at (M) {$M(x,y)$};
      \draw (O) -- (M) node[midway, above] {$R$};
    \end{tikzpicture}
    \caption*{Рисунок 9.10}
  \end{minipage}%
  \begin{minipage}{0.55\textwidth}
    $$
    x^2+y^2=R^2
    $$
    --- уравнение окружности с центром в начале координат
    (каноническое уравнение).
  \end{minipage}
\end{figure}

\textbf{Определение.} Эллипсом называется геометрическое место точек
на плоскости, сумма расстояний от которых до двух данных точек,
называемых фокусами, есть величина постоянная.

$$
|\overline{F_1M}|+|\overline{F_2M}|=2a,
$$
где $\overline{F_1M}$ и $\overline{F_2M}$ --- фокальные радиусы.

\begin{equation}
  \sqrt{(x+c)^2+y^2}+\sqrt{(x-c)^2+y^2}
  =
  2a. \tag{9.15}
\end{equation}

Полагая $b^2=a^2-c^2$, из (9.15) нетрудно получить уравнение
\begin{equation}
  \frac{x^2}{a^2}+\frac{y^2}{b^2}=1. \tag{9.16}
\end{equation}

\begin{figure}[H]
  \centering
  \begin{minipage}{0.4\textwidth}
    \centering
    \begin{tikzpicture}[scale=0.9, >=Latex]
      \draw[->] (-3.6,0) -- (3.8,0) node[right] {$x$};
      \draw[->] (0,-2.5) -- (0,2.7) node[above] {$y$};
      \node[below left] at (0,0) {$0$};
      \draw[thick] (0,0) ellipse (3 and 1.8);
      \coordinate (F1) at (-1.8,0);
      \coordinate (F2) at (1.8,0);
      \coordinate (M) at (1.0,1.7);
      \fill (F1) circle (1.3pt);
      \fill (F2) circle (1.3pt);
      \fill (M) circle (1.3pt);
      \node[below] at (F1) {$F_1(-c,0)$};
      \node[below] at (F2) {$F_2(c,0)$};
      \node[above right] at (M) {$M(x,y)$};
      \node[below] at (-3,0) {$A_1$};
      \node[below] at (3,0) {$A_2$};
      \node[left] at (0,1.8) {$B_1$};
      \node[left] at (0,-1.8) {$B_2$};
      \node[below] at (0.9,0) {$c$};
      \node[below] at (2.4,0) {$a$};
      \node[left] at (0,0.9) {$b$};
    \end{tikzpicture}
    \caption*{Рисунок 9.11}
  \end{minipage}%
  \begin{minipage}{0.55\textwidth}
    (9.16) --- каноническое уравнение эллипса.

    $|\overline{F_1F_2}|$ --- фокусное расстояние; $a$ и $b$ ---
    большая и малая полуоси эллипса.

    $$
    \varepsilon=\frac ca
    $$
    --- эксцентриситет; $0<\varepsilon<1$,

    $\varepsilon$ --- мера «сплюснутости»; $\varepsilon=0$ --- окружность.
  \end{minipage}
\end{figure}

\textbf{Определение.} Гипербола --- это геометрическое место точек на
плоскости, модуль разности расстояний от которых до двух данных
точек, называемых фокусами, есть величина постоянная.

$$
\left||\overline{F_1M}|-|\overline{F_2M}|\right|=2a.
$$

\begin{equation}
  \left|
  \sqrt{(x+c)^2+y^2}
  -
  \sqrt{(x-c)^2+y^2}
  \right|
  =
  2a. \tag{9.17}
\end{equation}

Путем эквивалентных преобразований получим уравнение
\begin{equation}
  \frac{x^2}{a^2}-\frac{y^2}{b^2}=1. \tag{9.18}
\end{equation}

\begin{figure}[H]
  \centering
  \begin{minipage}{0.4\textwidth}
    \centering
    \begin{tikzpicture}[scale=0.75, >=Latex]
      \draw[->] (-4.4,0) -- (4.6,0) node[right] {$x$};
      \draw[->] (0,-3.0) -- (0,3.2) node[above] {$y$};
      \node[below left] at (0,0) {$0$};
      \draw[dashed] (-3.6,-2.4) -- (3.6,2.4);
      \draw[dashed] (-3.6,2.4) -- (3.6,-2.4);
      \draw[thick, domain=1:3.9, samples=100] plot (\x,{0.7*sqrt(\x*\x-1)});
      \draw[thick, domain=1:3.9, samples=100] plot (\x,{-0.7*sqrt(\x*\x-1)});
      \draw[thick, domain=-3.9:-1, samples=100] plot (\x,{0.7*sqrt(\x*\x-1)});
      \draw[thick, domain=-3.9:-1, samples=100] plot (\x,{-0.7*sqrt(\x*\x-1)});
      \fill (-2,0) circle (1.3pt);
      \fill (2,0) circle (1.3pt);
      \node[below] at (-2,0) {$F_1(-c,0)$};
      \node[below] at (2,0) {$F_2(c,0)$};
      \node[above] at (0.7,0) {$c$};
      \node[above] at (1,0) {$a$};
      \fill (2.5,1.15) circle (1.3pt);
      \node[above right] at (2.5,1.15) {$M(x,y)$};
    \end{tikzpicture}
    \caption*{Рисунок 9.12}
  \end{minipage}%
  \begin{minipage}{0.55\textwidth}
    (9.18) --- каноническое уравнение гиперболы, где
    $$
    b^2=c^2-a^2.
    $$

    $a$ и $b$ --- действительная и мнимая полуоси гиперболы.

    $$
    \varepsilon=\frac ca
    $$
    --- эксцентриситет,
    $$
    x=\pm\frac a\varepsilon
    $$
    --- директрисы гиперболы и эллипса; прямые
    $$
    y=\pm\frac bax
    $$
    --- асимптоты гиперболы.
  \end{minipage}
\end{figure}

\textbf{Определение.} Парабола --- это геометрическое место точек на
плоскости, для которых расстояние до данной точки, называемой
фокусом, равно расстоянию до некоторой прямой, называемой директрисой.

\begin{figure}[H]
  \centering
  \begin{minipage}{0.4\textwidth}
    \centering
    \begin{tikzpicture}[scale=0.85, >=Latex]
      \draw[->] (-1.4,0) -- (4.5,0) node[right] {$x$};
      \draw[->] (0,-2.6) -- (0,2.8) node[above] {$y$};
      \node[below left] at (0,0) {$0$};
      \draw[thick, domain=0:3.4, samples=100] plot (\x,{sqrt(1.6*\x)});
      \draw[thick, domain=0:3.4, samples=100] plot (\x,{-sqrt(1.6*\x)});
      \draw[dashed] (-0.8,-2.2) -- (-0.8,2.2);
      \node[below] at (-0.8,0) {$-\frac p2$};
      \coordinate (F) at (0.8,0);
      \fill (F) circle (1.3pt);
      \node[below right] at (F) {$F\left(\frac p2,0\right)$};
      \coordinate (M) at (2.2,1.88);
      \coordinate (K) at (-0.8,1.88);
      \fill (M) circle (1.3pt);
      \fill (K) circle (1.3pt);
      \node[right] at (M) {$M(x,y)$};
      \node[left] at (K) {$K$};
      \draw (K) -- (M);
      \draw (F) -- (M);
    \end{tikzpicture}
    \caption*{Рисунок 9.13}
  \end{minipage}%
  \begin{minipage}{0.55\textwidth}
    $p$ --- параметр параболы $(p>0)$.

    $$
    KM=FM.
    $$

    $$
    x+\frac p2=
    \sqrt{\left(x-\frac p2\right)^2+y^2}.
    $$

    Путем эквивалентных преобразований получим уравнение
    \begin{equation}
      y^2=2px. \tag{9.19}
    \end{equation}

    (9.19) --- каноническое уравнение параболы.
  \end{minipage}
\end{figure}

\subsection*{\textbf{Оптические свойства кривых второго порядка}}

\begin{figure}[H]
  \centering
  \begin{minipage}{0.4\textwidth}
    \centering
    \begin{tikzpicture}[scale=0.85, >=Latex]
      \draw[->] (-3.8,0) -- (3.8,0) node[right] {$x$};
      \draw[->] (0,-2.3) -- (0,2.5) node[above] {$y$};
      \node[below left] at (0,0) {$0$};
      \draw[thick] (0,0) ellipse (3 and 1.7);
      \coordinate (F1) at (-1.6,0);
      \coordinate (F2) at (1.6,0);
      \coordinate (M) at (0.9,1.62);
      \fill (F1) circle (1.3pt);
      \fill (F2) circle (1.3pt);
      \fill (M) circle (1.3pt);
      \node[below] at (F1) {$F_1$};
      \node[below] at (F2) {$F_2$};
      \node[above right] at (M) {$M$};
      \draw[->] (F1) -- (M);
      \draw[->] (M) -- (F2);
    \end{tikzpicture}
    \caption*{Рисунок 9.14}
  \end{minipage}%
  \begin{minipage}{0.55\textwidth}
    Луч, выходящий из одного фокуса эллипса, после отражения от
    эллипса проходит через другой фокус. Для параболы лучи,
    параллельные оси, после отражения проходят через фокус.
  \end{minipage}
\end{figure}

\subsection*{\textbf{Полярные координаты}}

$\varphi$ --- полярный угол.

$$
\varphi\in[0,2\pi]\quad \text{или}\quad \varphi\in[-\pi,\pi].
$$

\begin{figure}[H]
  \centering
  \begin{minipage}{0.4\textwidth}
    \centering
    \begin{tikzpicture}[scale=0.95, >=Latex]
      \draw[->] (-0.2,0) -- (4.4,0) node[right] {$x$};
      \draw[->] (0,-0.2) -- (0,3.4) node[above] {$y$};
      \node[below left] at (0,0) {$0$};
      \coordinate (M) at (3.0,2.0);
      \draw[->, thick] (0,0) -- (M) node[midway, above] {$r$};
      \fill (M) circle (1.3pt);
      \node[right] at (M) {$M(r,\varphi)$};
      \draw (1.0,0) arc (0:33.7:1.0);
      \node at (1.25,0.35) {$\varphi$};
    \end{tikzpicture}
    \caption*{Рисунок 9.15}
  \end{minipage}%
  \begin{minipage}{0.55\textwidth}
    Связь декартовых координат с полярными:
    $$
    \begin{cases}
      x=r\cos\varphi,\\
      y=r\sin\varphi.
    \end{cases}
    $$
  \end{minipage}
\end{figure}

Получим уравнение окружности в полярных координатах:
$$
x^2+y^2=R^2,
$$
если $R=a$, где $a$ --- заданное число, то
$$
r^2(\cos^2\varphi+\sin^2\varphi)=a^2
\Rightarrow r=a.
$$

\begin{equation}
  \begin{cases}
    x=a\cos\varphi,\\
    y=a\sin\varphi.
  \end{cases}
  \tag{9.20}
\end{equation}

\begin{equation}
  \begin{cases}
    x=a\cos\varphi,\\
    y=b\sin\varphi.
  \end{cases}
  \tag{9.21}
\end{equation}

(9.20) --- параметрическое уравнение окружности.

(9.21) --- параметрическое уравнение эллипса.

\subsection*{\textbf{Пример}}

Составить уравнение параллели и перпендикуляра к прямой $L:2x-y+3=0$,
проходящих через точку $M(1;2)$.

\subsection*{\textbf{Решение}}

Точка $M\notin L$. Пусть $L_1\parallel L$; $L_2\perp L$;

т. $M\in L_1$; т. $M\in L_2$.

$$
A(x-x_0)+B(y-y_0)=0
$$
--- уравнение прямой, проходящей через т. $M(x_0;y_0)$ с нормальным
вектором $\vn(A,B)$.

$$
\Rightarrow L_1:2(x-1)-1(y-2)=0
\Rightarrow
$$
$$
\Rightarrow 2x-y=0
$$
--- уравнение $L_1$.

\begin{figure}[H]
  \centering
  \begin{minipage}{0.4\textwidth}
    \centering
    \begin{tikzpicture}[scale=0.75, >=Latex]
      \draw[->] (-2.5,0) -- (4.5,0) node[right] {$x$};
      \draw[->] (0,-2.5) -- (0,4.0) node[above] {$y$};
      \node[below left] at (0,0) {$0$};
      \draw[thick] (-1.2,0.6) -- (2.4,3.0) node[right] {$L_1$};
      \draw[thick] (-1.5,3.5) -- (3.5,0.2) node[right] {$L_2$};
      \draw[dashed] (-1.0,-1.0) -- (3.0,1.6) node[right] {$L$};
      \fill (1,2) circle (1.5pt);
      \node[above] at (1,2) {$M$};
    \end{tikzpicture}
    \caption*{Рисунок 9.16}
  \end{minipage}%
  \begin{minipage}{0.55\textwidth}
    На рисунке показаны прямая $L$ и проходящие через точку $M$
    параллельная ей прямая $L_1$ и перпендикулярная ей прямая $L_2$.
  \end{minipage}
\end{figure}

\subsection*{\textbf{Пример}}

В $\triangle ABC$ найти уравнение биссектрисы угла $B$, если
$$
A(3;-5);\qquad B(1;-3);\qquad C(2;-2).
$$

\subsection*{\textbf{Решение}}

$$
\frac{AB}{BC}=\frac{AH}{HC}.
$$

$$
AB=\sqrt{(1-3)^2+(-3+5)^2}=2\sqrt2.
$$

$$
BC=\sqrt{(2-1)^2+(-2+3)^2}=\sqrt2.
$$

$$
\frac{AB}{BC}=\frac{2\sqrt2}{\sqrt2}=2;
\qquad
x=\frac{x_1+\lambda x_2}{1+\lambda},
\qquad
y=\frac{y_1+\lambda y_2}{1+\lambda}.
$$

$$
\Rightarrow x=\frac{3+2\cdot 2}{1+2}=\frac73;
$$
$$
y=\frac{-5+2\cdot(-2)}{1+2}=-3.
$$

\begin{figure}[H]
  \centering
  \begin{minipage}{0.4\textwidth}
    \centering
    \begin{tikzpicture}[scale=0.85, >=Latex]
      \coordinate (A) at (0,0);
      \coordinate (B) at (2.2,1.5);
      \coordinate (C) at (4.1,0.2);
      \coordinate (H) at (2.9,0.1);
      \draw[thick] (A) -- (B) -- (C) -- cycle;
      \draw[thick] (B) -- (H);
      \fill (A) circle (1.3pt) node[left] {$A$};
      \fill (B) circle (1.3pt) node[above] {$B$};
      \fill (C) circle (1.3pt) node[right] {$C$};
      \fill (H) circle (1.3pt) node[below] {$H$};
    \end{tikzpicture}
    \caption*{Рисунок 9.17}
  \end{minipage}%
  \begin{minipage}{0.55\textwidth}
    $$
    \overline{BH}\left(\frac73-1;-3-(-3)\right)
    =
    \left(\frac43;0\right)
    $$
    $$
    \Rightarrow \vq(1,0)
    $$
    --- направляющий вектор для прямой $BH$.
  \end{minipage}
\end{figure}

$$
\frac{x-1}{1}=\frac{y+3}{0}
\Rightarrow
0\cdot(x-1)=1\cdot(y+3).
$$

$$
y=-3
$$
--- уравнение биссектрисы $BH$.

\subsection*{\textbf{Пример}}

Найти расстояние от точки $M(3;4)$ до прямой $L:x-2y+1=0$.

\subsection*{\textbf{Решение}}

$$
d=
\left|
\frac{Ax_0+By_0+C}{\sqrt{A^2+B^2}}
\right|
=
\left|
\frac{1\cdot3-2\cdot4+1}{\sqrt{1^2+(-2)^2}}
\right|
=
\frac4{\sqrt5}.
$$

\subsection*{\textbf{Пример}}

Найти угол между прямыми $L_1:2x-y+3=0$ и $L_2:2x+3y+5=0$.

\subsection*{\textbf{Решение}}

$$
\bar n_1(2,-1);\qquad \bar n_2(2;3).
$$

Обозначим
$$
\alpha=(\hat L_1,L_2).
$$

$$
\cos\alpha=
\left|
\frac{(\bar n_1,\bar n_2)}{|\bar n_1||\bar n_2|}
\right|
=
\left|
\frac{2\cdot2-1\cdot3}{\sqrt{4+1}\sqrt{4+9}}
\right|
=
\frac1{\sqrt{65}}.
$$

$$
\alpha=\arccos\frac1{\sqrt{65}}.
$$

\subsection*{\textbf{Пример}}

Установить, какую кривую определяет уравнение
$$
y^2+4=x^2.
$$
Сделать чертеж.

\subsection*{\textbf{Решение}}

$$
x^2-y^2=4.
$$

$$
\frac{x^2}{4}-\frac{y^2}{4}=1
$$
--- гипербола.

$$
a=2;\qquad b=2;\qquad b^2=c^2-a^2.
$$

$$
c^2=a^2+b^2=4+4=8;\qquad c=2\sqrt2.
$$

$$
y=\pm\frac bax=\pm x
$$
--- асимптоты.

$$
F_1(-2\sqrt2;0);\qquad F_2(2\sqrt2;0)
$$
--- фокусы.

$$
A_1(-2,0);\qquad A_2(2,0)
$$
--- вершины гиперболы.

$$
B_1(0,-2);\qquad B_2(0,2)
$$
--- мнимые вершины.

\begin{figure}[H]
  \centering
  \begin{minipage}{0.4\textwidth}
    \centering
    \begin{tikzpicture}[scale=0.75, >=Latex]
      \draw[->] (-4.3,0) -- (4.4,0) node[right] {$x$};
      \draw[->] (0,-3.2) -- (0,3.4) node[above] {$y$};
      \node[below left] at (0,0) {$0$};
      \foreach \x in {-1,1} {
        \draw (\x,0.08) -- (\x,-0.08) node[below] {$\x$};
      }
      \foreach \y in {-1,1} {
        \draw (0.08,\y) -- (-0.08,\y) node[left] {$\y$};
      }
      \draw[dashed] (-3.6,-3.6) -- (3.6,3.6);
      \draw[dashed] (-3.6,3.6) -- (3.6,-3.6);
      \draw[thick, domain=2:4, samples=100] plot (\x,{sqrt(\x*\x-4)});
      \draw[thick, domain=2:4, samples=100] plot (\x,{-sqrt(\x*\x-4)});
      \draw[thick, domain=-4:-2, samples=100] plot (\x,{sqrt(\x*\x-4)});
      \draw[thick, domain=-4:-2, samples=100] plot (\x,{-sqrt(\x*\x-4)});
      \fill (-2.828,0) circle (1.3pt) node[below] {$F_1$};
      \fill (2.828,0) circle (1.3pt) node[below] {$F_2$};
      \fill (-2,0) circle (1.3pt) node[above] {$A_1$};
      \fill (2,0) circle (1.3pt) node[above] {$A_2$};
      \fill (0,2) circle (1.3pt) node[right] {$B_2$};
      \fill (0,-2) circle (1.3pt) node[right] {$B_1$};
    \end{tikzpicture}
    \caption*{Рисунок 9.18}
  \end{minipage}%
  \begin{minipage}{0.55\textwidth}
    Гипербола
    $$
    \frac{x^2}{4}-\frac{y^2}{4}=1
    $$
    с асимптотами $y=\pm x$, фокусами $F_1(-2\sqrt2;0)$,
    $F_2(2\sqrt2;0)$, вершинами $A_1(-2,0)$, $A_2(2,0)$ и мнимыми
    вершинами $B_1(0,-2)$, $B_2(0,2)$.
  \end{minipage}
\end{figure}

\subsection*{\textbf{Пример}}

Установить, какую кривую определяет уравнение
$$
4x^2+9y^2+16x-18y-11=0.
$$
Сделать чертеж.

\subsection*{\textbf{Решение}}

$$
4x^2+9y^2+16x-18y-11=0.
$$

$$
4(x^2+4x)+9(y^2-2y)-11=0.
$$

$$
4(x^2+4x+4-4)+
$$
$$
+9(y^2-2y+1-1)-11=
$$
$$
=4(x+2)^2+9(y-1)^2=36.
$$

$$
\frac{(x+2)^2}{9}+\frac{(y-1)^2}{4}=1
$$
--- эллипс.

т. $O_1(-2;1)$ --- центр эллипса.

$$
a=3;\qquad b=2.
$$

$$
b^2=a^2-c^2
\Rightarrow
$$
$$
\Rightarrow c^2=a^2-b^2=9-4=5.
$$

$$
c=\sqrt5.
$$

\begin{figure}[H]
  \centering
  \begin{minipage}{0.4\textwidth}
    \centering
    \begin{tikzpicture}[scale=0.65, >=Latex]
      \draw[->] (-6.2,0) -- (2.2,0) node[right] {$x$};
      \draw[->] (0,-2.0) -- (0,4.6) node[above] {$y$};
      \node[below left] at (0,0) {$O$};
      \draw[dashed, ->] (-2,1) -- (1.8,1) node[right] {$x'$};
      \draw[dashed, ->] (-2,1) -- (-2,4.3) node[above] {$y'$};
      \draw[thick] (-2,1) ellipse (3 and 2);
      \fill (-2,1) circle (1.3pt) node[below right] {$O_1$};
      \fill ({-2-sqrt(5)},1) circle (1.3pt) node[below] {$F_1$};
      \fill ({-2+sqrt(5)},1) circle (1.3pt) node[below] {$F_2$};
      \node[below] at (-5,0) {$-5$};
      \node[below] at (-2,0) {$-2$};
      \node[below] at (-1,0) {$-1$};
      \node[left] at (0,1) {$1$};
      \node[below] at (-5,1) {$-5$};
      \node[below] at (-2,1) {$O_1$};
      \node[above] at (-2,3) {$b$};
      \node[below] at (-0.5,1) {$a$};
    \end{tikzpicture}
    \caption*{Рисунок 9.19}
  \end{minipage}%
  \begin{minipage}{0.55\textwidth}
    $$
    F_1(-2-\sqrt5;1);\qquad F_2(-2+\sqrt5;1).
    $$

    $F_1$, $F_2$ --- фокусы эллипса.
  \end{minipage}
\end{figure}

\end{document}
